\subsection{Модуль lab4}

\subsubsection{Kirchhoff: число остовных деревьев}

\textbf{Вход:} \texttt{Graph<int> const\&}~-- неориентированный граф.

\textbf{Выход:} \texttt{long long}~-- число различных остовных деревьев.

\textbf{Описание:} Метод \texttt{count} строит матрицу Лапласа $B$
размера $n \times n$: $B[i][j] = -1$ для смежных вершин $i \neq j$,
$B[i][i]$~-- степень вершины $i$. Смежность определяется как
$\texttt{adj}[i][j] \neq 0$ или $\texttt{adj}[j][i] \neq 0$.

Затем берётся алгебраическое дополнение $B'$~-- матрица $B$ без нулевой
строки и нулевого столбца размера $(n-1) \times (n-1)$.

\begin{lstlisting}[style=cstyle, caption={Метод Гаусса с целочисленной арифметикой (kirchhoff.cpp)}]
for (int k = 0; k < m; k++) {
    int pivot = -1;
    for (int i = k; i < m; i++)
        if (mat[i][k] != 0) { pivot = i; break; }
    if (pivot == -1) return 0;
    if (pivot != k) swap(mat[k], mat[pivot]);

    for (int i = k + 1; i < m; i++) {
        for (int j = k + 1; j < m; j++) {
            mat[i][j] = mat[k][k] * mat[i][j]
                      - mat[i][k] * mat[k][j];
            if (k > 0)
                mat[i][j] /= mat[k-1][k-1];
        }
        mat[i][k] = 0;
    }
}
return abs(mat[m-1][m-1]);
\end{lstlisting}

\textbf{Сложность:} $O(V^3)$ по времени; $O(V^2)$ по памяти.

\subsubsection{Kruskal: алгоритм Краскала}

\textbf{Вход:} \texttt{Graph<int> const\&}~-- связный неориентированный
взвешенный граф.

\textbf{Выход:} \texttt{vector<Edge>}~-- рёбра минимального остова;
пустой вектор, если граф несвязен.

\textbf{Описание:} Конструктор собирает все рёбра (каждое один раз,
условие $i < j$) в вектор \texttt{edges}; веса берутся по модулю.
Метод \texttt{compute()} сортирует рёбра по весу и добавляет их в остов,
если они не создают цикла. Цикл проверяется через BFS по уже построенному
остову: если вершины $u$ и $v$ уже связаны~-- ребро пропускается.

\begin{lstlisting}[style=cstyle, caption={Основной цикл алгоритма Краскала (kruskal.cpp)}]
vector<Edge> Kruskal::compute() {
    sort(edges.begin(), edges.end());
    vector<Edge> T;
    Graph<int> tree(n);
    for (auto& e : edges) {
        if ((int)T.size() == n - 1) break;
        if (hasPath(e.u, e.v, tree)) continue;
        T.push_back(e);
        tree.addEdge(e.u, e.v, e.w);
        tree.addEdge(e.v, e.u, e.w);
    }
    return ((int)T.size() == n - 1) ? T : vector<Edge>{};
}
\end{lstlisting}

\textbf{Сложность:} $O(E^2)$ по времени (BFS на каждое ребро);
$O(V + E)$ по памяти.

\subsubsection{Prufer: код Прюфера}

\textbf{Вход (encode):} \texttt{vector<Edge> const\&}~-- рёбра МОД,
\texttt{int n}~-- число вершин.

\textbf{Выход (encode):} \texttt{PruferCode}~-- поля: \texttt{seq}
(последовательность длины $n-2$), \texttt{weights} (веса рёбер),
\texttt{n}.

\textbf{Описание (encode):} Строится временный граф из рёбер МОД.
Хранится вектор степеней \texttt{d} и вектор активных вершин \texttt{V}.
На каждом шаге ищется активная висячая вершина с минимальным номером
($d[v] = 1$); её единственный сосед записывается в последовательность,
вес ребра~-- в \texttt{weights}. Вершина $v$ удаляется.

\begin{lstlisting}[style=cstyle, caption={Кодирование в код Прюфера с сохранением весов (prufer.cpp)}]
PruferCode Prufer::encode(const vector<Edge>& mst, int p) {
    for (int i = 0; i < p - 1; i++) {
        int v = -1;
        for (int k = 0; k < p; k++)
            if (V[k] && d[k] == 1) { v = k; break; }
        int neighbor = -1, weight = 0;
        for (int to = 0; to < p; to++)
            if (V[to] && adj[v][to] != 0)
                { neighbor = to; weight = wgt[v][to]; break; }
        code.seq.push_back(neighbor);
        code.weights.push_back(weight);
        V[v] = false; d[v]--; d[neighbor]--;
    }
    return code;
}
\end{lstlisting}

\textbf{Вход (decode):} \texttt{PruferCode}~-- структура с полями
\texttt{seq}, \texttt{weights}, \texttt{n}.

\textbf{Выход (decode):} \texttt{vector<Edge>}~-- рёбра восстановленного
дерева с весами.

\textbf{Описание (decode):} Поддерживается вектор активных вершин \texttt{B}.
Для каждого элемента \texttt{seq[i]} ищется наименьшая неудаленная вершина вершина,
которой нет в остатке кода \texttt{seq[i..n-2]}.
Между ней и \texttt{seq[i]} добавляется ребро с весом \texttt{weights[i]},
вершина удаляется.

\begin{lstlisting}[style=cstyle, caption={Декодирование кода Прюфера (prufer.cpp)}]
vector<Edge> Prufer::decode(const PruferCode& code) {
    int p = code.n;
    vector<Edge> E;
    vector<bool> B(p, true);

    for (int i = 0; i < p - 1; i++) {
        int v = -1;
        for (int k = 0; k < p; k++) {
            if (B[k]) {
                bool in_rest = false;
                for (int j = i; j < p - 1; j++)
                    if (code.seq[j] == k) { in_rest = true; break; }
                if (!in_rest) { v = k; break; }
            }
        }
        E.push_back({v, code.seq[i], code.weights[i]});
        B[v] = false;
    }
    return E;
}
\end{lstlisting}

\textbf{Сложность:} $O(V^2)$ по времени; $O(V)$ по памяти.

\subsubsection{MIS: максимальное независимое множество}

\textbf{Вход:} \texttt{vector<vector<int>> const\&}~-- матрица смежности
исходного графа или МОД (по выбору пользователя).

\textbf{Выход:} \texttt{vector<int>}~-- вершины максимального
независимого множества.

\textbf{Описание:} Метод \texttt{compute()} запускает рекурсивный
\texttt{backtrack(S, T)}, где $S$~-- текущее независимое множество,
$T$~-- кандидаты. Для каждой вершины $v \in T$ проверяется, не смежна
ли она с вершинами из $S$. Если нет~-- добавляем $v$ и убираем из
кандидатов всех соседей $v$. Если новое множество больше текущего
максимума~-- сохраняем.

\begin{lstlisting}[style=cstyle, caption={Backtracking для нахождения MIS (mis.cpp)}]
void MIS::backtrack(vector<int> S, vector<int> T) {
    for (int v : T) {
        if (isIndependent(S, v)) {
            vector<int> S_plus_v = S;
            S_plus_v.push_back(v);
            if ((int)S_plus_v.size() > m) {
                X = S_plus_v;
                m = S_plus_v.size();
            }
            vector<int> newT;
            for (int u : T)
                if (u != v && adj[v][u] == 0 && adj[u][v] == 0)
                    newT.push_back(u);
            backtrack(S_plus_v, newT);
        }
    }
}
\end{lstlisting}

\textbf{Сложность:} $O(2^V)$ по времени в худшем случае; $O(V)$ по памяти.
